documentclass[11pt]{article}
    \title{\textbf{Gummi 0.8.3}}
    \author{Alexander van der Meij}
    \date{}
    
    \addtolength{\topmargin}{-3cm}
    \addtolength{\textheight}{3cm}
\begin{document}

\maketitle
\thispagestyle{empty}

\section{Introduction}
\subsection{Эффективный 1D-гамильтониан и параметрический резонанс}

Рассмотрим топологическую трубку солитона как замкнутую упругую струну, параметризованную безразмерной координатой $s \in [0, 2\pi)$. Вследствие симметрии узла-трилистника, самодействие витков (crosstalk) индуцирует эффективный периодический потенциал (искривление) с пространственной частотой $k=3$:
\begin{equation}
    V_{\text{cross}}(s) = V_0 \cos(3s)
\end{equation}
где $V_0$ — амплитуда геометрического напряжения в местах перекрестий узла.

Возбужденное состояние (нейтрон) моделируется как фазовая волна (фрейминг), распространяющаяся вдоль трубки. Из численного эксперимента установлено, что амплитуда волны равна $\sqrt{\alpha}$, а пространственная частота соответствует тау-моде $k=15$. Уравнение волны имеет вид:
\begin{equation}
    \psi(s) = \sqrt{\alpha} \cos(15s)
\end{equation}

Дефект масс $\Delta m$ вычисляется как энергия этой волны на фоне периодического потенциала узла:
\begin{equation}
    \Delta E = \int_0^{2\pi} \left[ \frac{1}{2} \left( \frac{\partial \psi}{\partial s} \right)^2 + \frac{1}{2} V_{\text{cross}}(s) \psi^2(s) \right] ds
\end{equation}

Подставляя $\psi(s)$ в кинетический член, получаем строгую аналитическую зависимость упругой энергии от постоянной тонкой структуры и квадрата топологического индекса:
\begin{equation}
    E_{\text{elastic}} = \frac{1}{2} \int_0^{2\pi} \left( -15 \sqrt{\alpha} \sin(15s) \right)^2 ds = \frac{225}{2} \alpha \pi
\end{equation}

Член взаимодействия (интерференция волны с геометрией узла) определяет стабильность частицы:
\begin{equation}
    E_{\text{int}} = \frac{V_0 \alpha}{2} \int_0^{2\pi} \cos(3s) \cos^2(15s) ds
\end{equation}
Поскольку частота тау-моды ($k=15$) строго кратна симметрии трилистника ($k=3$), интеграл перекрытия обеспечивает параметрическую стабилизацию (lock-in эффект), предотвращая мгновенный распад нейтрона и обеспечивая его аномально большое время жизни.
